\begin{tikzpicture}[x=1.00cm,y=0.64cm]
\path[use as bounding box] (0,-0.10) rectangle (16,33.75);

% --------------------------------------------------------------------------
% t = 3.
\begin{scope}[shift={(0,23.60)}]
\draw[wqstage] (0,0) rectangle (16,10.10);
\node[wqwordlabel,anchor=west] at (0.65,9.10) {$1122112$};

\coordinate (d10pos) at (1.50,2.10);
\coordinate (d11pos) at (5.00,2.10);
\coordinate (d12pos) at (13.00,2.10);
\coordinate (d20pos) at (5.50,0.75);
\coordinate (d21pos) at (10.00,0.75);
\coordinate (d22pos) at (15.00,0.75);
\coordinate (dcutlevel) at (0,3.65);
\coordinate (d1endone) at (3.00,3.65);
\coordinate (d2endone) at (7.00,3.65);
\coordinate (d1endtwo) at (11.00,3.65);
\coordinate (d2endtwo) at (14.00,3.65);

\coordinate (dvonepos) at (3.00,7.65);
\coordinate (dvtwopos) at (7.00,6.45);
\coordinate (dvthreepos) at (11.00,5.35);
\draw[wqweaveone] (2.00,8.80)--(dvonepos);
\draw[wqweaveone] (4.00,8.80)--(dvonepos);
\draw[wqweaveone] (dvonepos)--(d1endone);
\draw[wqweavetwo] (6.00,8.80)--(dvtwopos);
\draw[wqweavetwo] (8.00,8.80)--(dvtwopos);
\draw[wqweavetwo] (dvtwopos)--(d2endone);
\draw[wqweaveone] (10.00,8.80)--(dvthreepos);
\draw[wqweaveone] (12.00,8.80)--(dvthreepos);
\draw[wqweaveone] (dvthreepos)--(d1endtwo);
\draw[wqweavetwo] (14.00,8.80)--(d2endtwo);
\node[wqweavefrozen] (dvone) at (dvonepos) {};
\node[wqweavefrozen] (dvtwo) at (dvtwopos) {};
\node[wqweavefrozen] (dvthree) at (dvthreepos) {};
\draw[wqcut] (2.60,3.65)--(14.40,3.65);
\node[wqwordlabel,anchor=west] at (0.65,3.65) {$1212$};

\draw[wqqbox] (0.45,0.02) rectangle (15.55,3.15);
\node[wqqboxtitle] at (0.75,3.15)
  {$Q(1212)$};
\node[wqstringfrozen] (d10) at (d10pos) {$[0,1]$};
\node[wqpivot]        (d11) at (d11pos) {$[1,3]$};
\node[wqstringfrozen] (d20) at (d20pos) {$[0,2]$};
\node[wqstringmutable](d21) at (d21pos) {$[2,4]$};
\node[wqstringfrozen] (d12) at (d12pos) {$[3,\infty]$};
\node[wqstringfrozen] (d22) at (d22pos) {$[4,\infty]$};
% All edge families are superimposed as continuous arrows.
% Three opposite 1/2 pairs: the complete v1--v2--v3 block.
  \draw[wqcancelblue] (dvtwo) to[bend left=5] (dvone);
  \draw[wqcancelorange] (dvone) to[bend left=5] (dvtwo);
  \draw[wqcancelblue] (dvthree) to[bend left=5] (dvtwo);
  \draw[wqcancelorange] (dvtwo) to[bend left=5] (dvthree);
  \draw[wqcancelblue] (dvone) to[bend left=16] (dvthree);
  \draw[wqcancelorange] (dvthree) to[bend right=22] (dvone);
% Terminal-weight and string-quiver edges.
  \wqextensiontriangle{d10}{dvone}{d11}
  \wqextensiontriangle{d20}{dvtwo}{d21}
  \wqextensiontriangle{d11}{dvthree}{d12}
  \draw[wqstringarrow] (d11)--(d10);
  \draw[wqstringhalf] (d10)--(d20);
  \draw[wqstringarrow] (d12)--(d11);
  \draw[wqstringarrow] (d20)--(d11);
  \draw[wqstringarrow] (d11)--(d21);
  \draw[wqstringarrow] (d21)--(d12);
  \draw[wqstringhalf] (d12)--(d22);
  \draw[wqstringarrow] (d21)--(d20);
  \draw[wqstringarrow] (d22)--(d21);

% Edge labels are drawn after the edge system.
\path (dvtwo) to[bend left=5]
  node[wqhalfblue,pos=0.44,below=0.22mm] {$\frac12$} (dvone);
\path (dvone) to[bend left=5]
  node[wqhalforange,pos=0.56,above=0.22mm] {$\frac12$} (dvtwo);
\path (dvthree) to[bend left=5]
  node[wqhalfblue,pos=0.44,below=0.22mm] {$\frac12$} (dvtwo);
\path (dvtwo) to[bend left=5]
  node[wqhalforange,pos=0.56,above=0.22mm] {$\frac12$} (dvthree);
\path (dvone) to[bend left=16]
  node[wqhalfblue,pos=0.60,below=0.32mm] {$\frac12$} (dvthree);
\path (dvthree) to[bend right=22]
  node[wqhalforange,pos=0.40,above=0.32mm] {$\frac12$} (dvone);
\path (d10)--node[wqhalf,pos=0.52] {$\frac12$} (d20);
\path (d12)--node[wqhalf,pos=0.50] {$\frac12$} (d22);
\node[wqvlabel,left=0.8mm of dvone] {$\sigma_1$};
\node[wqvlabel,above=0.8mm of dvtwo] {$\sigma_2$};
\node[wqvlabel,right=0.8mm of dvthree] {$\sigma_3$};
\end{scope}

% Fourth mutation.
\draw[wqtimearrow] (8.00,23.42)--node[wqmovelabel,right]
  {$\mu_{(1,[1,3])}\quad 121\to212$} (8.00,22.08);

% --------------------------------------------------------------------------
% t = 4.
\begin{scope}[shift={(0,11.80)}]
\draw[wqstage] (0,0) rectangle (16,10.10);
\node[wqwordlabel,anchor=west] at (0.65,9.10) {$1122112$};

\coordinate (e20pos) at (2.50,0.75);
\coordinate (e21pos) at (6.00,0.75);
\coordinate (e22pos) at (12.00,0.75);
\coordinate (e23pos) at (15.00,0.75);
\coordinate (e10pos) at (5.50,2.10);
\coordinate (e11pos) at (8.50,2.10);
\coordinate (ecutlevel) at (0,3.65);
\coordinate (e2endone) at (4.00,3.65);
\coordinate (e1endone) at (7.00,3.65);
\coordinate (e2endtwo) at (10.50,3.65);
\coordinate (e2endthree) at (14.00,3.65);

\coordinate (evonepos) at (3.00,7.65);
\coordinate (evtwopos) at (7.00,6.45);
\coordinate (evthreepos) at (11.00,5.35);
\coordinate (ebraid) at (7.00,4.85);
\draw[wqweaveone] (2.00,8.80)--(evonepos);
\draw[wqweaveone] (4.00,8.80)--(evonepos);
\draw[wqweaveone] (evonepos)--(ebraid);
\draw[wqweavetwo] (6.00,8.80)--(evtwopos);
\draw[wqweavetwo] (8.00,8.80)--(evtwopos);
\draw[wqweavetwo] (evtwopos)--(ebraid);
\draw[wqweaveone] (10.00,8.80)--(evthreepos);
\draw[wqweaveone] (12.00,8.80)--(evthreepos);
\draw[wqweaveone] (evthreepos)--(ebraid);
\draw[wqweavetwo] (14.00,8.80)--(e2endthree);
\draw[wqweavetwo] (ebraid)--(e2endone);
\draw[wqweaveone] (ebraid)--(e1endone);
\draw[wqweavetwo] (ebraid)--(e2endtwo);
\node[wqweavefrozen] (evone) at (evonepos) {};
\node[wqweavefrozen] (evtwo) at (evtwopos) {};
\node[wqweavefrozen] (evthree) at (evthreepos) {};
\node[wqbraid] at (ebraid) {$\times$};
\draw[wqcut] (3.60,3.65)--(14.40,3.65);
\node[wqwordlabel,anchor=west] at (0.65,3.65) {$2122$};

\draw[wqqbox] (0.45,0.02) rectangle (15.55,3.15);
\node[wqqboxtitle] at (0.75,3.15)
  {$Q(2122)$};
\node[wqstringfrozen] (e20) at (e20pos) {$[0,1]$};
\node[wqstringfrozen] (e10) at (e10pos) {$[0,2]$};
\node[wqstringmutable](e21) at (e21pos) {$[1,3]$};
\node[wqstringfrozen] (e11) at (e11pos) {$[2,\infty]$};
\node[wqpivot]        (e22) at (e22pos) {$[3,4]$};
\node[wqstringfrozen] (e23) at (e23pos) {$[4,\infty]$};
% Frozen--frozen arrows.
% v1 -> v3 is 1/2 + 1/2; the adjacent pairs cancel.
  \draw[wqcancelblue] (evtwo) to[bend left=5] (evone);
  \draw[wqcancelorange] (evone) to[bend left=5] (evtwo);
  \draw[wqcancelblue] (evthree) to[bend left=5] (evtwo);
  \draw[wqcancelorange] (evtwo) to[bend left=5] (evthree);
  \draw[wqsharedarrow] (evone)
    .. controls (5.40,8.75) and (9.15,8.15) .. (evthree);
% Terminal-weight and string-quiver arrows are superimposed continuously.
  \wqextensiontriangle{e20}{evtwo}{e22}
  \wqextensiontriangle{e20}{evthree}{e21}
  \wqextensiontriangle
    [out=165,in=-125,looseness=1.05]
    [out=20,in=75,looseness=1.55]
    {e21}{evone}{e22}
  \draw[wqstringarrow] (e11)--(e10);
  \draw[wqstringhalf] (e20)--(e10);
  \draw[wqstringhalf] (e11)--(e23);
% Remaining string-quiver arrows.
\draw[wqstringarrow] (e10)--(e21);
\draw[wqstringarrow] (e21)--(e11);
\draw[wqstringarrow] (e21)--(e20);
\draw[wqstringarrow] (e22)--(e21);
\draw[wqstringarrow] (e23)--(e22);

% Edge labels are drawn after the edge system.
\path (evone) .. controls (5.40,8.75) and (9.15,8.15) ..
  node[wqhalfsum,pos=0.52,above=0.35mm]
  {$\textcolor{wqorange}{\frac12}+\textcolor{wqblue}{\frac12}$}
  (evthree);
\path (evtwo) to[bend left=5]
  node[wqhalfblue,pos=0.44,below=0.22mm] {$\frac12$} (evone);
\path (evone) to[bend left=5]
  node[wqhalforange,pos=0.56,above=0.22mm] {$\frac12$} (evtwo);
\path (evthree) to[bend left=5]
  node[wqhalfblue,pos=0.44,below=0.22mm] {$\frac12$} (evtwo);
\path (evtwo) to[bend left=5]
  node[wqhalforange,pos=0.56,above=0.22mm] {$\frac12$} (evthree);
\path (e20)--node[wqhalf,pos=0.50] {$\frac12$} (e10);
\path (e11)--node[wqhalf,pos=0.32] {$\frac12$} (e23);
\node[wqvlabel,left=0.8mm of evone] {$\sigma_1$};
\node[wqvlabel,above=0.8mm of evtwo] {$\sigma_2$};
\node[wqvlabel,right=0.8mm of evthree] {$\sigma_3$};
\end{scope}

% Fifth mutation.
\draw[wqtimearrow] (8.00,11.62)--node[wqmovelabel,right]
  {$\mu_{(2,[3,4])}\quad 22\to2$} (8.00,10.28);

% --------------------------------------------------------------------------
% t = 5.  This is the supplied terminal composite picture, with the same
% three exchange blocks and with Q(212) directly below the complete weave.
\begin{scope}
\draw[wqstage] (0,0) rectangle (16,10.10);
\node[wqwordlabel,anchor=west] at (0.65,9.10) {$1122112$};

\coordinate (f20pos) at (2.50,0.75);
\coordinate (f21pos) at (7.00,0.75);
\coordinate (f22pos) at (15.00,0.75);
\coordinate (f10pos) at (5.50,2.10);
\coordinate (f11pos) at (8.50,2.10);
\coordinate (fcutlevel) at (0,3.65);
\coordinate (f2endone) at (4.00,3.65);
\coordinate (f1endone) at (7.00,3.65);
\coordinate (f2endtwo) at (12.25,3.65);

\coordinate (fvonepos) at (3.00,7.95);
\coordinate (fvtwopos) at (7.00,7.05);
\coordinate (fvthreepos) at (11.00,6.15);
\coordinate (fbraid) at (7.00,5.35);
\coordinate (fvfourpos) at (12.25,4.55);
\draw[wqweaveone] (2.00,8.80)--(fvonepos);
\draw[wqweaveone] (4.00,8.80)--(fvonepos);
\draw[wqweaveone] (fvonepos)--(fbraid);
\draw[wqweavetwo] (6.00,8.80)--(fvtwopos);
\draw[wqweavetwo] (8.00,8.80)--(fvtwopos);
\draw[wqweavetwo] (fvtwopos)--(fbraid);
\draw[wqweaveone] (10.00,8.80)--(fvthreepos);
\draw[wqweaveone] (12.00,8.80)--(fvthreepos);
\draw[wqweaveone] (fvthreepos)--(fbraid);
\draw[wqweavetwo] (14.00,8.80)--(fvfourpos);
\draw[wqweavetwo] (fbraid)--(f2endone);
\draw[wqweaveone] (fbraid)--(f1endone);
\draw[wqweavetwo] (fbraid)--(fvfourpos);
\draw[wqweavetwo] (fvfourpos)--(f2endtwo);
\node[wqcontraction] (fvone) at (fvonepos) {};
\node[wqweavefrozen] (fvtwo) at (fvtwopos) {};
\node[wqweavefrozen] (fvthree) at (fvthreepos) {};
\node[wqbraid] at (fbraid) {$\times$};
\node[wqweavefrozen] (fvfour) at (fvfourpos) {};
\draw[wqcut] (3.60,3.65)--(12.65,3.65);
\node[wqwordlabel,anchor=west] at (0.65,3.65) {$212$};

\draw[wqqbox] (0.45,0.02) rectangle (15.55,3.15);
\node[wqqboxtitle] at (0.75,3.15)
  {$Q(212)$};
\node[wqstringfrozen] (f20) at (f20pos) {$[0,1]$};
\node[wqstringfrozen] (f10) at (f10pos) {$[0,2]$};
\node[wqstringmutable](f21) at (f21pos) {$[1,3]$};
\node[wqstringfrozen] (f11) at (f11pos) {$[2,\infty]$};
\node[wqstringfrozen] (f22) at (f22pos) {$[3,\infty]$};
% All colored edge families are superimposed as continuous arrows.
% Frozen--frozen block: blue is the intrinsic Q_W contribution.  On
% v4 -> v2 the terminal K-pullback contributes another 1/2 in orange, so the
% coincident stroke displays the sum 1.  On v3--v4 the two shallow opposing
% arcs expose canceling contributions.
  \draw[wqcancelblue] (fvthree) to[bend left=6] (fvfour);
  \draw[wqcancelorange] (fvfour) to[bend left=6] (fvthree);
  \draw[wqsharedarrow] (fvfour)--(fvtwo);
  \draw[wqlocalarrow] (fvone)
    .. controls (5.45,8.80) and (9.15,8.15) .. (fvthree);
  \draw[wqlocalarrow] (fvfour)
    .. controls (12.80,9.10) and (5.50,9.45) .. (fvone);
% C_beta: terminal-weight arrows end at the actual weave vertices.
  \wqextensiontriangle{f20}{fvtwo}{f21}
  \wqextensiontriangle{f20}{fvthree}{f21}
  \wqextensiontriangle[out=13,in=-141,looseness=0.9]
    {f21}{fvfour}{f22}
% String-quiver edges.
  \draw[wqstringarrow] (f11)--(f10);
  \draw[wqstringhalf] (f20)--(f10);
  \draw[wqstringarrow] (f10)--(f21);
  \draw[wqstringarrow] (f21)--(f11);
  \draw[wqstringhalf] (f11)--(f22);
  \draw[wqstringarrow] (f21)--(f20);
  \draw[wqstringarrow] (f22)--(f21);

% Edge labels are drawn after the edge system.
\path (fvfour)--node[wqhalfsum,pos=0.48,above=0.32mm]
  {$\textcolor{wqorange}{\frac12}+\textcolor{wqblue}{\frac12}$}
  (fvtwo);
\path (fvthree) to[bend left=6]
  node[wqhalfblue,pos=0.44,above=0.20mm] {$\frac12$} (fvfour);
\path (fvfour) to[bend left=6]
  node[wqhalforange,pos=0.56,below=0.20mm] {$\frac12$} (fvthree);
\path (f20)--node[wqhalf,pos=0.48] {$\frac12$} (f10);
\path (f11)--node[wqhalf,pos=0.50] {$\frac12$} (f22);
\node[wqvlabel,left=0.8mm of fvone] {$\sigma_1$};
\node[wqvlabel,above=0.8mm of fvtwo] {$\sigma_2$};
\node[wqvlabel,right=0.8mm of fvthree] {$\sigma_3$};
\node[wqvlabel,right=0.8mm of fvfour] {$\sigma_4$};
\end{scope}

\end{tikzpicture}
